The double copy is a map between gauge, gravity, and scalar theories that arose out of the study of relationships
between the scattering amplitudes of those theories. It has implications for the study of classical solutions,
which it relates via various generally compatible formalisms. This dissertation presents research that investigates
various extensions of the classical double copy. An extension of the Weyl double copy into five dimensions is
presented, along with a prototypical example from a large class of highly symmetric spacetimes. The extension is
shown to be consistent with earlier results on the Kerr-Schild double copy of higher dimensional spacetimes. Next,
a proposal is made for the double copy structure of black hole horizons. Under the double copy, horizons seem to
map to quantities that are computable in the corresponding gauge and scalar theories, but resist physical
interpretation. The last extension explored is of expressing the double copy structure of generic spacetimes by
zooming into a null geodesic, which corresponds to taking the Penrose limit, and then expanding about the geodesic
order by order in a scaling parameter. Known instances of double copied spacetimes are shown to be consistent with
the Penrose limit and expansions about it. For arbitrary spacetimes, the double copy relations are shown to
generically fail at the second order in the expansion about a null geodesic.
